% =====================================================================
%  Deliverable 1 — Generative Artificial Intelligence (580694)
%  Universidad de Concepción · 31 de agosto de 2026
%  Joaquín Sepúlveda · Renato Saavedra · Lucas Saldivia
%
%  UNA PÁGINA, compilado desde LaTeX. Overleaf → pdfLaTeX.
% =====================================================================
\documentclass[9pt,a4paper,landscape]{extarticle}

\usepackage[T1]{fontenc}
\usepackage[utf8]{inputenc}
\usepackage[spanish,es-noquoting]{babel}
% babel-spanish vuelve activa la comilla doble: "c se convierte en ç y rompe
% el JSON de ejemplo. Esto la desactiva.
\shorthandoff{"}
\usepackage{libertine}
\usepackage[scaled=0.80]{beramono}
\usepackage{microtype}
\usepackage[margin=7mm]{geometry}
\usepackage{multicol}
\usepackage{xcolor}
\usepackage{booktabs}
\usepackage{tikz}
\usepackage{array}
\usepackage{colortbl}

\definecolor{tinta}{HTML}{12191A}
\definecolor{acento}{HTML}{B8501C}
\definecolor{verde}{HTML}{146B5F}
\definecolor{ambar}{HTML}{C08A24}
\definecolor{gris}{HTML}{5A6A6B}
\definecolor{fondo}{HTML}{EDF1F0}
\definecolor{linea}{HTML}{C8D4D1}

\color{tinta}
\pagestyle{empty}
\setlength{\columnsep}{5mm}
\setlength{\parindent}{0pt}
\setlength{\parskip}{0.85mm}
\renewcommand{\arraystretch}{0.98}

\newcommand{\seccion}[1]{\vspace{1.0mm}{\bfseries\footnotesize\color{acento}\MakeUppercase{#1}}\par\nobreak\vspace{0.4mm}\nobreak}
\newcommand{\sub}[1]{{\bfseries #1}\par\nobreak\vspace{-1.3mm}\nobreak}
\newcommand{\q}{\textquotedbl}   % comilla recta, inmune a los shorthands de babel
\newcommand{\cod}[1]{{\ttfamily\small #1}}
\newcommand{\destaque}[1]{%
  \begingroup\setlength{\fboxsep}{2.2mm}%
  \colorbox{fondo}{\parbox{\dimexpr\columnwidth-4.4mm\relax}{#1}}%
  \endgroup\par}

\begin{document}
\raggedright

% ---------------------------------------------------------------- cabecera
{\LARGE\bfseries ¿Puede un modelo de menos de 8B validar una inscripción de ramos?}\\[1mm]
{\footnotesize\color{gris}%
\textbf{Deliverable 1} · Generative Artificial Intelligence (580694) · Universidad de Concepción · 31 de agosto de 2026\\
\textbf{Equipo:} Joaquín Sepúlveda · Renato Saavedra · Lucas Saldivia \quad\textbullet\quad
\textbf{Repositorio:} \texttt{github.com/Luquitas02/genai-inscripcion-ramos}}

\vspace{1mm}{\color{linea}\rule{\textwidth}{0.9pt}}

\begin{multicols}{3}

% =====================================================================
%  COLUMNA 1
% =====================================================================
\seccion{La tarea}

Dado el historial académico de un estudiante y una consulta escrita como la
escribiría una persona real, decidir si puede inscribir el ramo \textbf{y citar
la regla que sostiene la decisión}.

\destaque{\small\itshape
``oye, me falta Cálculo 2 pero la estoy tomando ahora, ¿podría inscribir
ecuaciones el que viene?''\par\vspace{1.2mm}
\normalfont\ttfamily\scriptsize
\{\q decision\q:\q condicional\q, \q regla\q:\q R-DEPENDE-APROBACION\q\}}

\vspace{0.5mm}
Ecuaciones Diferenciales exige Álgebra II, ya aprobada, y Cálculo II, que está
\emph{en curso}. Hoy no está aprobada y el próximo semestre sí. Por este
semestre la respuesta sería \cod{no}, con el mismo expediente.

\vspace{0.5mm}
\sub{Qué cuenta como correcto}
Coincidencia exacta de los dos campos contra un verificador determinista, y se
reportan por separado.

\vspace{0.5mm}
\sub{La decisión}
\cod{sí} puede inscribir · \cod{no} algo lo impide y no se resuelve solo ·
\cod{condicional} puede, sujeto a aprobar lo que cursa ahora \emph{o} a que
proceda la excepción por carga bajo el mínimo.

\vspace{0.5mm}
\sub{La regla}
El código de la asignatura que falta (\cod{527150}) o el artículo que decide:
\cod{R-DEPENDE-APROBACION}, \cod{R-EXCEPCION-PRERREQ},
\cod{R-CREDITOS-MINIMOS}, \cod{R-TOPE-MAX}, \cod{R-ESPECIAL-*},
\cod{R-SIN-IMPEDIMENTO}. Quince en los 60 casos: 8 códigos, 7 artículos.

\seccion{Dominio y datos}

Malla de Ingeniería Civil Industrial UdeC, congelada como fuente única para los
cuatro entregables. \textbf{61} asignaturas, 227 créditos, cadenas de
prerrequisitos de hasta \textbf{6} niveles, \textbf{14} ramos con umbral de
créditos aprobados y 4 reglas globales.

\vspace{0.5mm}
Auditada antes de usarse: sin ciclos, créditos consistentes, un umbral
inalcanzable detectado y corregido contra el portal.

\seccion{Cómo se mide}

\textbf{60 casos} sintéticos. La respuesta de cada uno la fija el verificador.
Decisiones balanceadas 21/21/18, así que \textbf{la mejor constante trivial
rinde 35{,}0\,\%}.

\vspace{0.5mm}
\sub{Generación hacia atrás}
Se elige primero qué regla debe decidir y se construye el historial que la
produce. El caso se descarta si el verificador no confirma lo declarado.

\vspace{0.5mm}
\sub{Control de sesgo}
Cruce rasgo $\times$ respuesta antes de medir: ningún rasgo superficial de la
pregunta predice la decisión por sí solo.

\vspace{0.5mm}
\sub{Tres condiciones}
Zero-shot en prosa, zero-shot con la pregunta reducida a código y período, y
few-shot con tres ejemplos resueltos.

% =====================================================================
%  COLUMNA 2
% =====================================================================
\seccion{Cada modelo colapsa en una categoría distinta}

Reparto de las 60 respuestas de cada modelo, en zero-shot con prosa:

\vspace{1mm}
\begin{tikzpicture}[x=1mm,y=1mm]
  \def\H{5.2}
  \def\G{10.5}
  \def\U{1.0333}

  % --- referencia ARRIBA: hace de clave de colores y de patron a comparar ---
  \fill[verde!35]  (0,{3*\G})       rectangle ++({18*\U},\H);
  \fill[acento!35] ({18*\U},{3*\G}) rectangle ++({21*\U},\H);
  \fill[ambar!35]  ({39*\U},{3*\G}) rectangle ++({21*\U},\H);
  \draw[gris,line width=0.4pt] (0,{3*\G}) rectangle ({60*\U},{3*\G+\H});
  \node[font=\tiny\bfseries,tinta] at ({9*\U},{3*\G+\H/2})    {18 sí};
  \node[font=\tiny\bfseries,tinta] at ({28.5*\U},{3*\G+\H/2}) {21 no};
  \node[font=\tiny\bfseries,tinta] at ({49.5*\U},{3*\G+\H/2}) {21 condicional};
  \node[font=\scriptsize\bfseries,anchor=west] at (0,{3*\G+\H+2.2}) {lo correcto};

  % --- Qwen: 9 / 48 / 3
  \fill[verde]  (0,{2*\G})       rectangle ++({9*\U},\H);
  \fill[acento] ({9*\U},{2*\G})  rectangle ++({48*\U},\H);
  \fill[ambar]  ({57*\U},{2*\G}) rectangle ++({3*\U},\H);
  \node[font=\tiny\bfseries,white] at ({4.5*\U},{2*\G+\H/2})  {9};
  \node[font=\tiny\bfseries,white] at ({33*\U},{2*\G+\H/2})   {48};
  \node[font=\tiny\bfseries,white] at ({58.5*\U},{2*\G+\H/2}) {3};
  \node[font=\scriptsize\bfseries,anchor=west] at (0,{2*\G+\H+2.2}) {Qwen2.5-7B \textcolor{gris}{\normalfont 7{,}6\,B}};

  % --- Phi: 20 / 8 / 31, mas 1 salida ilegible
  \fill[verde]   (0,{1*\G})       rectangle ++({20*\U},\H);
  \fill[acento]  ({20*\U},{1*\G}) rectangle ++({8*\U},\H);
  \fill[ambar]   ({28*\U},{1*\G}) rectangle ++({31*\U},\H);
  \fill[gris!30] ({59*\U},{1*\G}) rectangle ++({1*\U},\H);
  \node[font=\tiny,gris,anchor=west] at ({60*\U+1.4},{1*\G+\H/2}) {1 sin leer};
  \node[font=\tiny\bfseries,white] at ({10*\U},{1*\G+\H/2})   {20};
  \node[font=\tiny\bfseries,white] at ({24*\U},{1*\G+\H/2})   {8};
  \node[font=\tiny\bfseries,white] at ({43.5*\U},{1*\G+\H/2}) {31};
  \node[font=\scriptsize\bfseries,anchor=west] at (0,{1*\G+\H+2.2}) {Phi-3.5-mini \textcolor{gris}{\normalfont 3{,}8\,B}};

  % --- Mistral: 50 / 6 / 4
  \fill[verde]  (0,0)       rectangle ++({50*\U},\H);
  \fill[acento] ({50*\U},0) rectangle ++({6*\U},\H);
  \fill[ambar]  ({56*\U},0) rectangle ++({4*\U},\H);
  \node[font=\tiny\bfseries,white] at ({25*\U},{\H/2}) {50};
  \node[font=\tiny\bfseries,white] at ({53*\U},{\H/2}) {6};
  \node[font=\tiny\bfseries,white] at ({58*\U},{\H/2}) {4};
  \node[font=\scriptsize\bfseries,anchor=west] at (0,{\H+2.2}) {Mistral-7B \textcolor{gris}{\normalfont 7{,}2\,B}};
\end{tikzpicture}

\vspace{1mm}
\textbf{Ninguno deduce la respuesta del expediente.} Cada uno se vuelca sobre
una categoría distinta.

\seccion{Resultados}

{\scriptsize
\begin{tabular}{@{}llrrr@{}}
\toprule
\textbf{Modelo} & \textbf{Condición} & \textbf{Dec.} & \textbf{Regla} & \textbf{Ambas} \\
\midrule
Mistral 7{,}2B & zero, prosa & 40{,}0 & 35{,}0 & 33{,}3 \\
           & zero, reducida & 36{,}7 & 31{,}7 & 31{,}7 \\
           & few-shot & 41{,}7 & \textbf{45{,}0} & \textbf{38{,}3} \\
\addlinespace[0.4mm]
Qwen 7{,}6B & zero, prosa & 40{,}0 & 8{,}3 & 8{,}3 \\
           & zero, reducida & 38{,}3 & 6{,}7 & 6{,}7 \\
           & few-shot & 43{,}3 & 31{,}7 & 31{,}7 \\
\addlinespace[0.4mm]
Phi 3{,}8B & zero, prosa & 43{,}3 & 6{,}7 & 5{,}0 \\
           & zero, reducida & 35{,}0 & 5{,}0 & 5{,}0 \\
           & few-shot & 35{,}0 & 23{,}3 & 16{,}7 \\
\midrule
\rowcolor{fondo}
\multicolumn{2}{@{}l}{\itshape Constante trivial} & \itshape 35{,}0 & \itshape 30{,}0 & \itshape 30{,}0 \\
\bottomrule
\end{tabular}}

\vspace{0.8mm}
{\scriptsize\color{gris}Porcentajes sobre 60 casos, generación determinista. La
constante trivial responde siempre lo mismo sin leer nada: \cod{no} para la
decisión, \cod{sí} con \cod{R-SIN-IMPEDIMENTO} para el par.}

\vspace{0.8mm}
\textbf{Cinco de las nueve condiciones quedan por debajo de esa constante} en
la métrica conjunta. La mejor le gana por 8{,}3 puntos.

\seccion{Factibilidad de ejecución}

{\scriptsize
\begin{tabular}{@{}lr@{}}
Hardware & GPU T4 gratuita (Colab), cuantización 4 bits \\
VRAM usada & $\approx$5\,GB de 15 disponibles \\
Tiempo por caso & 4{,}2\,s (Phi) -- 11{,}0\,s (Mistral) \\
Tokens de entrada & 3\,515 -- 3\,914 (Mistral) \\
Corrida completa & 9 condiciones, $\approx$66\,min de inferencia \\
Salidas con JSON válido & \textbf{539 / 540} \\
\end{tabular}}

\vspace{0.6mm}
El formato casi nunca falló. Los errores fueron de razonamiento.

% =====================================================================
%  COLUMNA 3
% =====================================================================
\seccion{Diagnóstico de la falla}

\sub{E1 · No representa ``cursándola ahora''}
De las 21 \cod{condicional} esperadas, Qwen produce 3 y Mistral 4; Phi produce
31. Los 18 casos con un prerrequisito en curso se parten 12 \cod{condicional} /
6 \cod{no} según el período preguntado. Mistral acierta \textbf{3 de esos 18}, y
con la pregunta reducida acierta \textbf{0}: responde \cod{sí} a los 18.

\vspace{0.6mm}
\sub{E2 · Depende de que la pregunta le apunte}
Al reducir la consulta a código y período (el historial completo sigue en el
prompt), el acierto \textbf{cae en los tres modelos}: $-1{,}7$, $-3{,}3$ y
$-8{,}3$ puntos.

\vspace{0.6mm}
\sub{E3 · No copia un identificador que tiene delante}
Phi-3.5-mini degrada la cadena \cod{R-CREDITOS-MINIMOS}, escrita literalmente
en su prompt, en formas inexistentes:

\vspace{-1mm}
{\ttfamily\scriptsize R-CREDITOS-MINIMISMU, R-CREDITOS-MINIMISIMO,
R-CREDITOS-MINIMISMUDA}

\vspace{0.6mm}
Más invenciones puras (\cod{R-8-CURSADO}, \cod{R-0321}).
\textbf{35 instancias en Phi; 0 en Qwen; 0 en Mistral.}

\vspace{0.6mm}
\sub{E4 · Se contradice dentro de su propia respuesta}
Responde \cod{sí} citando una regla que bloquea: hasta \textbf{28 de 60} en
Qwen con la pregunta reducida. Mistral lo hace 10 veces, siempre con
\cod{R-YA-CURSADA}, que no es correcta en ninguno de los 60 casos.

\vspace{0.8mm}
\destaque{\footnotesize
\textbf{Los ejemplos ordenan la salida pero no mejoran la decisión.}\par\vspace{0.6mm}
El few-shot baja las contradicciones (Qwen a 0, Mistral de 10 a 3, Phi de 18 a
11), casi elimina las alucinaciones (Phi de 35 a 4) y sube el acierto de regla
en los tres ($+23{,}4$, $+16{,}6$ y $+10{,}0$ puntos). \textbf{La decisión
apenas se mueve:} $+3{,}3$, $-8{,}3$ y $+1{,}7$.}

\seccion{Modelos candidatos}

{\scriptsize
\begin{tabular}{@{}p{0.35\columnwidth}rp{0.36\columnwidth}@{}}
\toprule
\textbf{Modelo} & \textbf{Par.} & \textbf{Rol} \\
\midrule
Mistral-7B-Instruct-v0.3 \textcolor{gris}{Mistral AI} & 7{,}2\,B & mejor acierto conjunto; base del harness \\
Qwen2.5-7B-Instruct \textcolor{gris}{Alibaba} & 7{,}6\,B & multilingüe; el que más gana con ejemplos \\
Phi-3.5-mini-instruct \textcolor{gris}{Microsoft} & \textbf{3{,}8\,B} & prueba de si la capacidad es necesaria \\
\bottomrule
\end{tabular}}

\vspace{0.8mm}
\textbf{Criterio de selección.} Tres familias distintas: tres tamaños del mismo
modelo no dirían nada sobre arquitectura ni sobre entrenamiento en español. Los
tres son open-weight, sin licencia restringida y corren en la T4 gratuita.

\vspace{0.8mm}
\textbf{Sobre el candidato bajo el techo.} Propusimos Phi-3.5-mini, con menos
de la mitad de los parámetros permitidos, bajo la hipótesis de que una salida
de veinte tokens no necesita capacidad sobrante. \textbf{La medición la
refuta.} Es el peor y el único que corrompe identificadores, aunque el más
rápido (\textbf{2{,}4$\times$} sobre Mistral).

\seccion{Cómo crece la tarea}

La dificultad es un eje medible, así que se endurece sin cambiar de tema.

\vspace{0.4mm}
{\scriptsize
\textbf{D2} prompt y contexto contra E1 y E2 \quad
\textbf{D3} consulta al grafo contra E3, validador contra E4 \quad
\textbf{D4} puntaje final contra este baseline.}

\end{multicols}
\end{document}
